\scanpage{117}
Comportement fonctoriel des groupes d'inertie.

\handwritten{1.} Soient $X$ un schéma, $U$ un ouvert de $X$, $\handwritten{G}$ un groupe profini, $U'$ un
revêtement principal de $U$ de groupe $G\deleted{\unsure{'}}$, $x$ un point \inserted{\handwritten{\unsure{géom|géom.}}} de $X$, $\bar X$ le localisé
strict de $X$ en $x$, $\bar U$ l'image inverse de $U$ dans $\bar X$, $\bar U'=U'\times_U\bar U$. On suppose
donné de plus un schéma $\bar X'$ sur $\bar X$, avec des opérations de $G$ sur $\bar X'$, de
telle façon que le \deleted{\unsure{espace|space|l'espace}}$\inserted{\bar U\text{-schéma}}$ à opérateurs $\bar X'\handwritten{\mid}\bar U$ co\deleted{\unsure{p}}ïncide avec $\bar U'$.

Pratiquement, on sera intéressé par le cas où $x\handwritten{\in}X-U$ et où $x$ est
adhérent à $U$, où $\bar X'$ est entier sur $\bar X$ et où $\underline O_{\bar X}=\bar p(\underline O_{\overline{X'}})^G$, $\bar p:\overline{X'}\handwritten{\to}\bar X$ étant
la projection canonique, - de sorte que $G$ opère transitiv\deleted{\unsure{i|x}}ement sur
l'ensembl\unsure{e|é} fibre $\overline{X'}_{\bar x}$. De plus, on \deleted{\unsure{se place|se plac}} sera \inserted{\handwritten{le plus souvent}} \deleted{\unsure{d}}sous des conditions où
on peut construire canoniquement $\bar X'$ en termes de $\bar U'$; par exemple, si on
suppose $\deleted{\unsure{O|U}}X$ normal en $x$, donc $\bar X\deleted{\unsure{'}}$ normal, et $\bar U$ non vide, on définira $\bar X'$ comme
le normalisé de $\bar U$ relativement à $\bar U'$ sur $\bar U$ (i.e. comme limite projective
des nor\deleted{\unsure{al|ma}}malisés d\deleted{\unsure{es}}ans les $\overline{U'(i)}$ sur $\bar U$, $U'$ \deleted{\unsure{u}}étant é\deleted{\unsure{e}}crit comme limite projective de schémas principaux $U'(i)$ sur les groupes quotients finis $G(i)$ de
$G$).

Si maintenant \deleted{\unsure{$x$ est}} $\handwritten{\xi}$ est un point de $\overline{X'}_{\bar x}$, on peut considérer le
groupe de stabilité $G_{\handwritten{\xi}}$ de $\handwritten{\xi}$, formé des $g\in G$ tels que $g.\handwritten{\xi}=\handwritten{\xi}$; ce groupe
prendra le nom de \uline{groupe d'inertie} \uline{relatif à} $\handwritten{\xi}$. Dans le cas envisagé plus
haut où les points de $\bar X'$ sur $\bar x$ sont tous conjugués sous $G$, on voit que
pour $\handwritten{\xi}$ variable\deleted{\unsure{s}}, les $G_{\handwritten{\xi}}$ parcourent une classe de sous-groupes conjugués
de $G$, qu'on appelera \inserted{(par abus de langage)} \uline{les sous-groupes d'inertie} \deleted{de} \uline{relatifs à} $x$. Noter
qu'en toute rigueur cette classe de sous-groupes \deleted{d'inertie} n'est pas
déterminée par la seule connaissance de $x\handwritten{\in}X$, mais dépend du choix du
prolongement de $\overline{U'}$ en $\overline{X'}$. Lorsque $X$ est normal en $x$ et que $x$ est spécialisation d'un point de $U$ (ce qui signifie que $x$ est adhérent à $U$ si $\mathrm{esp}\handwritten{(}X\handwritten{)}$
est localement noethérien), il sera sous-entendu, \deleted{\unsure{-}}sauf mention du contraire, que le choix en question est fait par normalisation, comme expliqué
plus haut.
% archive sheet number: 116
